La calidad de la educación superior en Colombia es un tema que ha cobrado cada vez más relevancia tanto para las instituciones educativas como para la comunidad académica, especialmente en lo relacionado con la evaluación del desarrollo de competencias durante la trayectoria formativa de los estudiantes. En este contexto, el Instituto Colombiano para la Evaluación de la Educación (ICFES) aplica dos pruebas estandarizadas de gran importancia: la prueba Saber 11, que se presenta al finalizar la educación media, y la prueba Saber Pro, aplicada al culminar la formación universitaria de pregrado.
Estas evaluaciones permiten analizar el desempeño de los estudiantes en distintos momentos de su proceso educativo y constituyen una valiosa fuente de información para comprender cómo evolucionan sus competencias a lo largo del tiempo. Sin embargo, a pesar del potencial que ofrecen estos datos, pocas instituciones de educación superior los han utilizado de manera sistemática para relacionar los resultados obtenidos al ingreso y al final de la formación universitaria, identificar tendencias de mejora o estancamiento y respaldar decisiones académicas con evidencia objetiva. La Universidad de San Buenaventura también enfrenta este desafío, pero al mismo tiempo cuenta con una oportunidad importante para aprovechar la información histórica disponible y transformarla en conocimiento útil para fortalecer sus procesos académicos y de mejora continua.
Con el propósito de aportar en este sentido, este trabajo desarrolla un sistema de información con capacidades predictivas y comparativas orientado a analizar la evolución de las competencias de los estudiantes entre las pruebas Saber 11 y Saber Pro. Para ello, se integran los resultados de ambas evaluaciones junto con variables contextuales y socioeconómicas relevantes, lo que permite construir modelos estadísticos y de aprendizaje automático capaces de estimar el desempeño esperado de los estudiantes y medir el valor agregado que la institución aporta durante su proceso de formación.
Desde el campo de la ingeniería de datos y bajo este argumento, este proyecto representa una aplicación concreta de técnicas de minería de datos educativos (\textit{Educational Data Mining}, EDM), análisis estadístico, visualización interactiva al ámbito de la gestión académica universitaria y la integración de modelos basados en sistemas inteligentes. El uso de algoritmos como la regresión lineal, los árboles de decisión y el método de los $k$ vecinos más cercanos (KNN) permitirá construir modelos robustos de predicción, mientras que los tableros de control (\textit{dashboards}) facilitarán la interpretación de resultados por parte de los tomadores de decisiones institucionales.
